\bookchapter{Capítulo 9: Ofimática y estudio en movilidad}

\chapterintro{Yo también pensaba que trabajar en el teléfono era “apañarse”, hasta que un viaje de negocios me dejó sin computadora y tuve que terminar de revisar una propuesta y firmar un contrato desde el teléfono. Sinceramente, el cuello de botella del trabajo móvil muchas veces no está en el teléfono, sino en que no conoces sus límites.}

\booksection{I. Tratamiento de documentos: editar y crear desde la palma de la mano}

Editar documentos en la pantalla del teléfono se siente distinto que en la computadora. El mayor problema es la pantalla pequeña, los dedos gruesos y la maquetación fácilmente desalineada. Pero hay técnicas que mejoran bastante la experiencia:

\begin{itemize}
  \item \textbf{WPS Office}: la mejor localización, totalmente compatible con los formatos de Microsoft (docx/xlsx/pptx), con un conjunto de herramientas PDF muy completo. Me gusta especialmente su función de "nota de voz": cuando revisas un contrato y no es cómodo escribir, basta con dictar unos comentarios al teléfono y los convierte automáticamente en anotaciones de texto dentro del documento. Además, WPS incorpora una enorme cantidad de plantillas gratuitas que ahorran mucho tiempo al hacer un currículum o una presentación de informe.
  \item \textbf{Microsoft 365}: si tu empresa usa Office original, es la mejor opción. La compatibilidad con la versión de escritorío es casi perfecta y no aparecen los ligeros desajustes de maquetación que a veces se ven al abrir con WPS ciertos diseños complejos. Integra OneDrive en profundidad: el documento que dejaste a medías en la computadora se abre en el teléfono en su versión más reciente.
  \item \textbf{Shimo Docs / Tencent Docs}: estos dos son los representantes de la "colaboración en línea". Adecuados para que varias personas editen un mismo documento a la vez, con los cambios de cada uno en tiempo real. He usado Shimo Docs para preparar una guía de viaje con un amigo: los dos rellenábamos la tabla a la vez, mucho más eficiente que hacerlo uno solo.
\end{itemize}

\textbf{Flujo de trabajo práctico}: mi costumbre actual es guardar los archivos importantes en OneDrive, hacer la edición diaria con WPS y abrir Shimo Docs para la colaboración en equipo. Desde que tengo los archivos en la nube, nunca más me he visto en el apuro de “olvidar el pendrive”.

\booksection{II. App de notas: de apunte rápido a sistema de conocimiento}

A la hora de tomar notas, mucha gente quiere montar de entrada un sistema muy complejo y no aguanta ni tres días. Mi consejo es empezar ligero y luego pasar a algo más serío: primero crea el hábito de “anotar” y luego piensa en “gestionar”.

\begin{itemize}
  \item \textbf{Apunte rápido — basta con lo integrado en el sistema}: Notas de iOS arranca rapidísimo y admite escanear documentos, escribir a mano y grabar audio. Las notas de las marcas Android tampoco se quedan atrás y se pueden abrir rápido desde la pantalla de bloqueo. Sirven para anotar ideas temporales, listas de compra o puntos de reuniones. La clave es que abren rápido: cuando se te ocurre algo, lo anotas sin esperar a que cargue la App.
  \item \textbf{Leer más tarde — recopilar información}: la función "Flotante" de WeChat es en realidad una herramienta oculta. Si ves un buen artículo y no tienes tiempo de leerlo, pulsa "Flotante" en la esquina superior derecha y se contraerá en un punto circular pegado al borde de la pantalla; puedes abrirlo cuando quieras. Es mucho más cómodo que reenviarlo al "Asistente de transferencia de archivos". Además, herramientas de "leer más tarde" como Cubox o Pocket permiten recopilar en cualquier plataforma lo que veas.
  \item \textbf{Montar un sistema — Notíon y Lark}: cuando sientas que las notas sueltas empiezan a ser difíciles de encontrar, es el momento de Notíon. Es muy flexible: puedes crear una "base de datos de notas de lectura" que se organiza automáticamente por título, autor y puntuación; o montar un "tablero de proyecto" para visualizar el progreso de las tareas. Las "tablas multidimensiónales" de Lark son similares y el acceso desde China es más rápido. Llevo más de tres años usando Notíon para gestionar todos mis materiales de escritura, notas de estudio y progreso de proyectos.
\end{itemize}

\textbf{Estrategia práctica de notas}: mi costumbre es un “trío”: ideas temporales en la nota del sistema (apertura instantánea), artículos para leer en la ventana flotante de WeChat (a mano) y conocimiento que requiere acumulación a largo plazo en Notíon (sistematizado). No intentes meterlo todo en una sola herramienta desde el principio; así es más fácil abandonar.

\booksection{III. Escaneo y OCR: traslada el mundo en papel al teléfono}

El escaneo con el teléfono, bien usado, ahorra muchos viajes a la copistería. La última vez que tramité un visado necesitaba versiones electrónicas de un montón de documentos y los hice todos con el teléfono.

\begin{itemize}
  \item \textbf{Usuaríos de iPhone}: la función de escaneo integrada en Notas de iOS ya es muy buena. Abre Notas toca el icono de la cámara elige "Escanear documento"; al apuntar al archivo reconoce automáticamente el borde y corrige la perspectiva. Tras la captura puedes ajustar el color: el modo "Blanco y negro" da el resultado más parecido a un escáner; "Escala de grises" es adecuado para documentos con imágenes.
  \item \textbf{En conclusión:}: muchas marcas incluyen un modo "Escaneo de documentos" en la cámara. Si no lo encuentras, descarga Microsoft Lens desde la tienda de aplicaciones (gratis e integrado con OneDrive).
  \item \textbf{Reconocimiento OCR de texto}: si necesitas extraer el texto del escaneo para editarlo, hace falta OCR. La función "Extraer texto" integrada en WeChat es la más sencilla: mantén pulsada la imagen y se usa; la precisión es bastante buena. Para procesar grandes cantidades, puedes usar la App Baimiao o CamScanner. Para el reconocimiento de chino, el motor OCR de Baidu funciona bastante bien; la función de imagen a texto de WPS también es muy precisa.
  \item \textbf{Consejos de escaneo}: escanea en un lugar bien iluminado, preferiblemente con una lámpara lateral, para no proyectar la sombra de la mano sobre el documento. Si el documento tiene las esquinas dobladas, alísalo con la mano antes de escanear. Los documentos de varias páginas se pueden escanear de forma continua y combinarlos en un solo PDF, sin guardarlos uno a uno.
\end{itemize}

\booksection{IV. Qué disco en la nube elegir: no hay uno mejor, solo el que mejor se adapte a ti}

Estuve dándole vueltas mucho tiempo a qué disco en la nube elegir y al final vi que depende de cuál sea tu “ecosistema principal”.

\begin{itemize}
  \item \textbf{Kit completo de Apple iCloud}: sincronización tan fluida que ni notas su presencia. Las fotos que tomas con el teléfono aparecen en la computadora a los pocos minutos. Además, Apple hace la protección de privacidad más estricta. Pero los inconvenientes son evidentes: compartir con usuarios de Android o Windows es engorroso, solo 5 GB gratis y ampliar cuesta dinero.
  \item \textbf{Windows como principal OneDrive}: si usas una computadora con Windows y Office 365, OneDrive es la pareja perfecta. Al guardar un documento en Word, se guarda automáticamente en la nube. 1 TB junto con la suscripción a Office sale muy rentable. La función de copia de seguridad automática de fotos en la App móvil también es muy útil.
  \item \textbf{Para compartir en China Baidu Netdisk}: sinceramente, la limitación de velocidad de descarga de Baidu Netdisk se critica siempre, pero su ecosistema de compartir no tiene rival. Gran parte de los materiales de estudio y recursos audiovisuales en la red china se comparten a través de Baidu Netdisk. Si solo descargas de vez en cuando archivos compartidos por otros, se puede usar sin membresía. Si eres un recolector intensivo de recursos, abre la membresía; el tiempo que ahorras vale más.
  \item \textbf{Aliyun Drive / Quark Drive}: en los últimos años han surgido con fuerza. Aliyun Drive no limita la velocidad de subida y bajada (por ahora); Quark Drive está fuertemente integrado con el navegador Quark, lo que permite guardar directamente en el disco al buscar materiales, muy cómodo.
\end{itemize}

\textbf{Mi combinación}: iCloud para fotos personales y archivos diarios (privacidad primero), Baidu Netdisk para recursos y recibir compartidos (ecosistema primero), OneDrive para documentos de trabajo (integración con Office primero). Tres discos, cada uno con su función, sin interferir entre sí.

\booksection{V. Un escenario real de ofimática móvil}

Te cuento una experiencia real propia. Una vez, en un viaje de negocios a otra ciudad, el cliente me pidió cambiar cláusulas del contrato de urgencia. Mi flujo fue:

\begin{enumerate}
  \item En el hotel recibí por WeChat los comentarios del cliente (mensaje de texto + una foto con anotaciones manuscritas).
  \item Abrí OneDrive para encontrar el contrato original, lo abrí y edité con WPS, modificando las cláusulas una por una según los mensajes de WeChat del cliente.
  \item Al terminar, lo exporté con la función "Exportar como PDF" de WPS y lo envié al grupo de WeChat del cliente para confirmación.
  \item El cliente planteó una segúnda ronda de comentarios; usé la herramienta de lápiz de WPS para rodear directamente en el PDF las posiciónes modificadas, junto con notas de voz explicando los criterios.
  \item Tras confirmar que estaba todo correcto, usé la función "Firma electrónica" de WPS para firmar en línea y envié al otro para que firmara también.
  \item Por último, archivé el PDF firmado en Baidu Netdisk y en iCloud, y anoté una entrada en el registro de trabajo de Notíon.
\end{enumerate}

Todo el proceso, de principio a fin, duró unos cuarenta minutos, sin encender la computadora. Sinceramente, hace tres años no me creía que el teléfono pudiera hacer todo esto, pero hoy en día es posible. La clave es familiarizarse con estas herramientas con antelación y no ponerse a descargar y registrar cuentas a toda prisa cuando llegue el momento.

\begin{tcbpractice}{Prueba: monta tu flujo de ofimática móvil}
\begin{enumerate}
  \item Elige tu App principal de documentos (WPS o Microsoft 365), inicia sesión y vincula el disco en la nube.
  \item Abre Notas o la nota del sistema, usa la función de escaneo para capturar un documento, expórtalo como PDF y envíatelo.
  \item Elige una App de disco en la nube y configúrala para la copia de seguridad automática de fotos (elige "Solo con Wi-Fi" para ahorrar datos).
  \item Crea una página en Notíon o Lark e intenta montar una plantilla sencilla de "registro de trabajo".
\end{enumerate}

\end{tcbpractice}
